\documentclass[10pt,a4paper]{article}
\usepackage[T1]{fontenc}
\usepackage[utf8]{inputenc}
\usepackage[czech]{babel}
\usepackage{lmodern}
\usepackage{microtype}
\usepackage{geometry}
\geometry{left=14mm,right=14mm,top=18mm,bottom=18mm}
\usepackage[table]{xcolor}
\usepackage{array,longtable,tabularx,booktabs}
\usepackage{fancyhdr,lastpage}
\usepackage{ragged2e}
\usepackage{hyperref}
\hypersetup{hidelinks,pdftitle={Rozpis výuky matematiky – 3. ročník SPŠ, 102 hodin podle tematických plánů a skutečné praxe}}
\setlength{\parindent}{0pt}
\setlength{\parskip}{5pt}
\setlength{\tabcolsep}{4pt}
\renewcommand{\arraystretch}{1.08}
\setlength{\headheight}{18pt}
\pagestyle{fancy}\fancyhf{}
\fancyhead[L]{\footnotesize Matematika -- rozpis výuky}
\fancyhead[R]{\footnotesize 3. ročník SPŠ}
\fancyfoot[L]{\footnotesize Rozpis hodin podle tematického plánu}
\fancyfoot[R]{\footnotesize \thepage/\pageref{LastPage}}
\renewcommand{\headrulewidth}{0.35pt}
\newcommand{\testr}[1]{\textbf{#1}}
\begin{document}
\begin{center}
{\LARGE\bfseries Rozpis výuky matematiky}\\[3mm]
{\large 3. ročník střední průmyslové školy}\\[3mm]
{\small Přehled plánované výuky podle tematických plánů SP\v{S} a reálných zápisů výuky}
\end{center}
\vspace{3mm}
\begin{tabularx}{\textwidth}{@{}Xrrrr@{}}\toprule
\textbf{Ročník} & \textbf{Hodin} & \textbf{Krátké testy} & \textbf{Souhrnné} & \textbf{Testů celkem}\\\midrule
3. ročník & 102 & 8 & 4 & 12\\\midrule
\textbf{Celkem} & \textbf{102} & \textbf{8} & \textbf{4} & \textbf{12}\\\bottomrule
\end{tabularx}
\vspace{4mm}
\section*{Charakter dokumentu}
\begin{itemize}\setlength{\itemsep}{2pt}
\item Dokument obsahuje rozpis výuky matematiky pro 3. ročník SP\v{S} v celkové dotaci \textbf{102 hodin}.
\item Rozpis je členěn po měsících a čtvrtletích. Každý očíslovaný řádek představuje \textbf{jednu plánovanou vyučovací hodinu}.
\item Zkratka \textbf{TP} odkazuje na odpovídající tematický okruh tematického plánu. U testů je označení uvedeno přímo v náplni hodiny.
\item V každém čtvrtletí jsou zařazeny \textbf{tři testy}: dva průběžné testy (zpravidla 15--20 minut) a jedna souhrnná písemná práce.
\end{itemize}

\section*{Z čeho rozpis vychází}
\begin{itemize}\setlength{\itemsep}{2pt}
\item Hodinová dotace vychází z tematického plánu SP\v{S} pro 3. ročník: \textbf{102 hodin}.
\item Obsah a pořadí témat zohledňují reálné zápisy výuky v dostupných třídách, zejména v 1.A (2024/25), 2.A (2025/26), 3.A (2025/26) a 3.D (2025/26).
\item Historické zápisy slouží jako podklad pro reálné pořadí témat a odhad časové náročnosti; nejde o doslovný přepis třídních knih.
\end{itemize}

\section*{Podkladové dokumenty}
\begin{itemize}\setlength{\itemsep}{2pt}
\item tematické plány SP\v{S},
\item zápisy výuky.
\end{itemize}

\section*{Odlišnosti oproti tematickým plánům}
\begin{itemize}\setlength{\itemsep}{2pt}
\item \textbf{3. ročník:} pořadí tematických celků vychází ze skutečné výuky ve třídách 3.A a 3.D; zohledněno je i odlišné tempo obou tříd.
\end{itemize}
\clearpage

\section*{3. ročník – komplexní čísla a analytická geometrie}
\textit{Dotace 102 hodin. Čtvrtletní celky a tři testy v každém čtvrtletí jsou součástí hodinové dotace.}
\textit{Historická orientace: 3.A a 3.D postupovaly stejně po hlavních tématech, ale měly odlišná čísla hodin. Plán nepřiřazuje konkrétní historická čísla plánovaným hodinám.}
\vspace{2mm}
\footnotesize
\begin{longtable}{@{}>{\RaggedLeft\arraybackslash}p{9mm}>{\RaggedRight\arraybackslash}p{15mm}>{\RaggedRight\arraybackslash}p{150mm}@{}}
\toprule \textbf{Hod.} & \textbf{TP} & \textbf{Náplň vyučovací hodiny} \\\midrule\endfirsthead
\toprule \textbf{Hod.} & \textbf{TP} & \textbf{Náplň vyučovací hodiny} \\\midrule\endhead
\bottomrule\endfoot
\addlinespace[1mm]\multicolumn{3}{@{}l}{\footnotesize\textbf{1. čtvrtletí} -- 2 průběžné testy + 1 souhrnná práce}\\[0.5mm]
\addlinespace[1.1mm]\multicolumn{2}{@{}l}{\textbf{Září}} & \multicolumn{1}{r@{}}{\textbf{9 hodin}}\\
\midrule
1 & 1.1–1.7 & Úvod do třetího ročníku, opakování druhého ročníku. \\
2 & 1.1–1.7 & Zavedení komplexního čísla, imaginární jednotka a algebraický tvar. \\
3 & 1.1–1.7 & Sčítání a odčítání komplexních čísel. \\
4 & 1.1–1.7 & Násobení komplexních čísel v algebraickém tvaru. \\
5 & 1.1–1.7 & Komplexně sdružené číslo; sdruženost a absolutní hodnota. \\
6 & 1.1–1.7 & Dělení komplexních čísel a usměrnění jmenovatele. \\
7 & 1.1–1.7 & Operace s komplexními čísly – procvičení. \\
8 & 1.1–1.7 & Gaussova rovina: obraz komplexního čísla, modul, argument. \\
9 & 1.1–1.7 & Převod algebraického tvaru na goniometrický. \\
\addlinespace[1.1mm]\multicolumn{2}{@{}l}{\textbf{Říjen}} & \multicolumn{1}{r@{}}{\textbf{12 hodin}}\\
\midrule
10 & T1 & \testr{Průběžný test 1 (15–20 min): Operace s komplexními čísly v algebraickém tvaru. Zbytek hodiny: krátký rozbor a navazující procvičování.} \\
11 & 1.1–1.7 & Goniometrický tvar komplexního čísla: argument a modul. \\
12 & 1.1–1.7 & Násobení a dělení čísel v goniometrickém tvaru. \\
13 & 1.1–1.7 & Moivreova věta: mocniny komplexních čísel. \\
14 & 1.1–1.7 & Odmocňování komplexních čísel – geometrický význam. \\
15 & 1.1–1.7 & Kvadratická rovnice v oboru komplexních čísel. \\
16 & 1.1–1.7 & Binomické rovnice: základní příklady. \\
17 & 1.1–1.7 & Komplexní čísla: převody tvarů, operace a kontrola. \\
18 & 1.1–1.7 & Opakování komplexních čísel před souhrnným testem. \\
19 & 2.1–2.2 & Analytická geometrie v rovině: kartézská soustava souřadnic. \\
20 & T2 & \testr{Průběžný test 2 (15–20 min): Komplexní čísla – všechny probírané tvary. Zbytek hodiny: krátký rozbor a navazující procvičování.} \\
21 & 2.1–2.2 & Vzdálenost dvou bodů v rovině; vzorec a Pythagorova věta. \\
\addlinespace[1.1mm]\multicolumn{2}{@{}l}{\textbf{Listopad}} & \multicolumn{1}{r@{}}{\textbf{7 hodin}}\\
\midrule
22 & 2.1–2.2 & Střed úsečky: souřadnice a slovní úlohy. \\
23 & 2.1–2.2 & Trojúhelník v souřadnicích: délky stran a těžnice. \\
24 & 2.1–2.2 & Vektor jako orientovaná úsečka a jeho souřadnice. \\
25 & 2.1–2.2 & Velikost vektoru a jednotkový vektor. \\
26 & 2.1–2.2 & Sčítání a odčítání vektorů; geometrický význam. \\
27 & P1 & \testr{Čtvrtletní písemná práce: Komplexní čísla a úvod analytické geometrie.} \\
28 & 2.2–2.3 & Násobení vektoru reálným číslem; rovnoběžnost vektorů. \\
\addlinespace[1mm]\multicolumn{3}{@{}l}{\footnotesize\textbf{2. čtvrtletí} -- 2 průběžné testy + 1 souhrnná práce}\\[0.5mm]
29 & 2.2–2.3 & Lineární kombinace vektorů a lineární závislost. \\
30 & 2.2–2.3 & Skalární součin: definice a výpočet ze souřadnic. \\
31 & 2.2–2.3 & Kolmost a odchylka vektorů pomocí skalárního součinu. \\
32 & 2.2–2.3 & Skalární součin: smíšené úlohy a ověřování. \\
33 & 2.2–2.3 & Vektorový součin vektorů – význam a výpočet. \\
\addlinespace[1.1mm]\multicolumn{2}{@{}l}{\textbf{Prosinec}} & \multicolumn{1}{r@{}}{\textbf{9 hodin}}\\
\midrule
34 & 2.2–2.3 & Geometrický význam vektorového součinu: kolmý vektor. \\
35 & 2.2–2.3 & Obsah rovnoběžníku a trojúhelníku z vektorového součinu. \\
36 & T3 & \testr{Průběžný test 3 (15–20 min): Skalární a vektorový součin. Zbytek hodiny: krátký rozbor a navazující procvičování.} \\
37 & 2.2–2.3 & Vektorové operace: postup výpočtu a kontrola. \\
38 & 2.2–2.3 & Skalární a vektorový součin: výběr správné operace. \\
39 & 2.2–2.3 & Opakování vektorů a součinů před testem. \\
40 & 2.2–2.3 & Samostatné procvičení analytických úloh s vektory. \\
41 & 2.4–2.8 & Parametrické vyjádření přímky: bod a směrový vektor. \\
42 & 2.4–2.8 & Rovnice přímky z dvojice bodů: nalezení směrového vektoru. \\
\addlinespace[1.1mm]\multicolumn{2}{@{}l}{\textbf{Leden}} & \multicolumn{1}{r@{}}{\textbf{12 hodin}}\\
\midrule
43 & 2.4–2.8 & Bod na přímce: dosazení do parametrického tvaru. \\
44 & 2.4–2.8 & Parametrický tvar přímky: graf a procvičování. \\
45 & T4 & \testr{Průběžný test 4 (15–20 min): Vektory a parametrické vyjádření přímky. Zbytek hodiny: krátký rozbor a navazující procvičování.} \\
46 & 2.4–2.8 & Obecná rovnice přímky: normálový vektor. \\
47 & 2.4–2.8 & Převod parametrického vyjádření na obecnou rovnici. \\
48 & 2.4–2.8 & Převod obecné rovnice na parametrické vyjádření. \\
49 & 2.4–2.8 & Obecná rovnice přímky: poloha bodu vůči přímce. \\
50 & 2.4–2.8 & Směrnicový tvar přímky a jeho geometrický význam. \\
51 & 2.4–2.8 & Převody mezi třemi rovnicemi přímky. \\
52 & 2.4–2.8 & Vzájemná poloha dvou přímek: rovnoběžnost, totožnost, různoběžnost. \\
53 & P2 & \testr{Pololetní písemná práce: Komplexní čísla, vektory a přímky v rovině.} \\
54 & 2.4–2.8 & Výpočet průsečíku přímek z jejich rovnic. \\
\addlinespace[1mm]\multicolumn{3}{@{}l}{\footnotesize\textbf{3. čtvrtletí} -- 2 průběžné testy + 1 souhrnná práce}\\[0.5mm]
\addlinespace[1.1mm]\multicolumn{2}{@{}l}{\textbf{Únor}} & \multicolumn{1}{r@{}}{\textbf{9 hodin}}\\
\midrule
55 & 2.4–2.8 & Kolmost přímek a normálové vektory. \\
56 & 2.4–2.8 & Odchylka dvou přímek ze směrových vektorů. \\
57 & 2.4–2.8 & Vzdálenost bodu od přímky: odvození a výpočet. \\
58 & 2.4–2.8 & Vzdálenost dvou rovnoběžných přímek. \\
59 & 2.4–2.8 & Přímky v rovině: slovní a technické úlohy. \\
60 & 2.4–2.8 & Analytická geometrie přímek: úlohy s parametrem. \\
61 & 2.4–2.8 & Opakování rovnic přímky a vzájemné polohy. \\
62 & 2.4–2.8 & Opakování odchylky a vzdálenosti v rovině. \\
63 & 3.1–3.4 & Analytická geometrie v prostoru: souřadnice bodu a vektoru. \\
\addlinespace[1.1mm]\multicolumn{2}{@{}l}{\textbf{Březen}} & \multicolumn{1}{r@{}}{\textbf{12 hodin}}\\
\midrule
64 & 3.1–3.4 & Vzdálenost bodů a velikost vektoru v prostoru. \\
65 & T5 & \testr{Průběžný test 5 (15–20 min): Rovnice přímek, odchylky a vzdálenosti. Zbytek hodiny: krátký rozbor a navazující procvičování.} \\
66 & 3.1–3.4 & Parametrické vyjádření přímky v prostoru. \\
67 & 3.1–3.4 & Vzájemná poloha přímek v prostoru: různé případy. \\
68 & 3.1–3.4 & Mimoběžky a jejich rozpoznání pomocí rovnic. \\
69 & 3.1–3.4 & Parametrické vyjádření roviny: bod a dva směrové vektory. \\
70 & 3.1–3.4 & Obecná rovnice roviny a normálový vektor. \\
71 & 3.1–3.4 & Vzájemná poloha přímky a roviny; průsečík. \\
72 & 3.1–3.4 & Vzájemná poloha dvou rovin; průsečnice. \\
73 & 3.1–3.4 & Vzdálenost bodu od roviny – vzorec a podmínky. \\
74 & 3.1–3.4 & Odchylka přímek v prostoru: směrové vektory. \\
75 & T6 & \testr{Průběžný test 6 (15–20 min): Analytická geometrie v prostoru. Zbytek hodiny: krátký rozbor a navazující procvičování.} \\
\addlinespace[1.1mm]\multicolumn{2}{@{}l}{\textbf{Duben}} & \multicolumn{1}{r@{}}{\textbf{3 hodiny}}\\
\midrule
76 & 3.1–3.4 & Odchylka přímky a roviny; odchylka rovin. \\
77 & 3.1–3.4 & Vzdálenosti v prostoru: procvičení typových úloh. \\
78 & P3 & \testr{Čtvrtletní písemná práce: Přímky a roviny v prostoru, vzdálenosti.} \\
\addlinespace[1mm]\multicolumn{3}{@{}l}{\footnotesize\textbf{4. čtvrtletí} -- 2 průběžné testy + 1 souhrnná práce}\\[0.5mm]
79 & 4.1–4.9 & Kuželosečky: geometrická definice a společné souvislosti. \\
80 & 4.1–4.9 & Kružnice: středová rovnice a základní vlastnosti. \\
81 & 4.1–4.9 & Kružnice: obecná rovnice a doplnění na čtverec. \\
82 & 4.1–4.9 & Kružnice a přímka: sečna, tečna, vnější přímka. \\
83 & 4.1–4.9 & Tečna ke kružnici v daném bodě; normála. \\
84 & 4.1–4.9 & Kružnice: úlohy s parametrem a polohou přímky. \\
\addlinespace[1.1mm]\multicolumn{2}{@{}l}{\textbf{Květen}} & \multicolumn{1}{r@{}}{\textbf{9 hodin}}\\
\midrule
85 & 4.1–4.9 & Elipsa: ohnisková definice a geometrická konstrukce. \\
86 & 4.1–4.9 & Středová rovnice elipsy; poloosy a výstřednost. \\
87 & 4.1–4.9 & Elipsa z rovnice: určení středu, poloos a ohnisek. \\
88 & T7 & \testr{Průběžný test 7 (15–20 min): Kružnice a elipsa. Zbytek hodiny: krátký rozbor a navazující procvičování.} \\
89 & 4.1–4.9 & Elipsa: sestavení rovnice ze zadaných údajů. \\
90 & 4.1–4.9 & Přímka a elipsa: počet průsečíků, dosazení. \\
91 & 4.1–4.9 & Kružnice a elipsa: souhrnné procvičení a porovnání. \\
92 & 4.1–4.9 & Parabola: definice pomocí ohniska a řídicí přímky. \\
93 & 4.1–4.9 & Parabola: vrcholový tvar rovnice a orientace osy. \\
\addlinespace[1.1mm]\multicolumn{2}{@{}l}{\textbf{Červen}} & \multicolumn{1}{r@{}}{\textbf{9 hodin}}\\
\midrule
94 & 4.1–4.9 & Parabola z rovnice: ohnisko, řídicí přímka a náčrt. \\
95 & 4.1–4.9 & Přímka a parabola: průsečíky a tečna. \\
96 & 4.1–4.9 & Hyperbola: ohnisková definice, vrcholy a asymptoty. \\
97 & 4.1–4.9 & Hyperbola: základní rovnice a geometrický náčrt. \\
98 & 4.1–4.9 & Hyperbola: určení středu a poloos z rovnice. \\
99 & 4.1–4.9 & Srovnání elipsy, paraboly a hyperboly; rozlišovací znaky. \\
100 & T8 & \testr{Průběžný test 8 (15–20 min): Kuželosečky: parabola a hyperbola. Zbytek hodiny: krátký rozbor a navazující procvičování.} \\
101 & 4.1–4.9 & Kuželosečky: smíšené úlohy a opakování. \\
102 & P4 & \testr{Pololetní písemná práce: Kuželosečky a opakování analytické geometrie.} \\
\end{longtable}
\normalsize
\end{document}
